% ===== VARIANT: QUANTITATIVE RISK ANALYST ===========================
% Angle: regulatory capital, PD/LGD/DRC, VaR/ES, stress testing and
% model validation. Strongest evidenced profile — regulatory work first.
% ===================================================================
% =====================================================================
%  Shared preamble for all CV variants.
%  Each variant does:   % =====================================================================
%  Shared preamble for all CV variants.
%  Each variant does:   % =====================================================================
%  Shared preamble for all CV variants.
%  Each variant does:   \input{cv-preamble.tex}
%  Do not put content here — only packages, formatting and macros.
% =====================================================================
\documentclass[letterpaper,11pt]{article}

\usepackage{latexsym}
\usepackage[empty]{fullpage}
\usepackage{titlesec}
\usepackage{marvosym}
\usepackage[usenames,dvipsnames]{color}
\usepackage{verbatim}
\usepackage{enumitem}
\usepackage[hidelinks]{hyperref}
\usepackage{fancyhdr}
\usepackage[english]{babel}
\usepackage{tabularx}
\usepackage{fontawesome5}
\usepackage{multicol}
\usepackage{graphicx}
\setlength{\multicolsep}{-3.0pt}
\setlength{\columnsep}{-1pt}
\input{glyphtounicode}
\usepackage{paracol}

\pagestyle{fancy}
\fancyhf{}
\fancyfoot{}
\renewcommand{\headrulewidth}{0pt}
\renewcommand{\footrulewidth}{0pt}

\addtolength{\oddsidemargin}{-0.5in}
\addtolength{\evensidemargin}{-0.5in}
\addtolength{\textwidth}{0.9in}
\addtolength{\topmargin}{-.3in}
\addtolength{\textheight}{0.9in}
\hypersetup{
    colorlinks=true,
    linkcolor=blue,
    urlcolor=blue,
    filecolor=blue,
    pdfpagemode=FullScreen
}
\urlstyle{same}

\raggedbottom
\raggedright
\setlength{\tabcolsep}{0in}

\titleformat{\section}{
  \vspace{-4pt}\scshape\raggedright\large\bfseries
}{}{0em}{}[\color{black}\titlerule \vspace{-5pt}]

\pdfgentounicode=1

\newcommand{\resumeItem}[1]{
  \item\small{{#1 \vspace{-2pt}}}
}

\newcommand{\resumeSubheading}[4]{
  \vspace{-2pt}\item
    \begin{tabular*}{1.0\textwidth}[t]{l@{\extracolsep{\fill}}r}
      \textbf{#1} & \textbf{\small #2} \\
      \textit{\small#3} & \textit{\small #4} \\
    \end{tabular*}\vspace{-7pt}
}

\newcommand{\resumeProjectHeading}[2]{
    \item
    \begin{tabular*}{0.97\textwidth}{l@{\extracolsep{\fill}}r}
      \small#1 & \textbf{\small #2}\\
    \end{tabular*}\vspace{-7pt}
}

\renewcommand\labelitemi{$\vcenter{\hbox{\tiny$\bullet$}}$}
\renewcommand\labelitemii{$\vcenter{\hbox{\tiny$\bullet$}}$}

\newcommand{\resumeSubHeadingListStart}{\begin{itemize}[leftmargin=0.0in, label={}]}
\newcommand{\resumeSubHeadingListEnd}{\end{itemize}}
\newcommand{\resumeItemListStart}{\begin{itemize}}
\newcommand{\resumeItemListEnd}{\end{itemize}\vspace{-5pt}}

% ---------------------------------------------------------------------
% \cvheader{positioning line}{focus line}
% ---------------------------------------------------------------------
\newcommand{\cvheader}[2]{%
\begin{center}
    {\Huge \scshape Franck Deba} \\[0.4em]
    {\small \textit{#1}} \\[0.25em]
    {\small #2}
\end{center}
\vspace{0.5pt}
\small
\raisebox{-0.1\height}{\faPhone}~+33 (0)6 05 77 47 51 \hspace{1em}
\raisebox{-0.2\height}{\faEnvelope}~\href{mailto:debafranck@gmail.com}{\underline{debafranck@gmail.com}} \hspace{1em}
\raisebox{-0.2\height}{\faLinkedin}~\href{https://linkedin.com/in/deba-franck}{\underline{deba-franck}} \hspace{1em}
\raisebox{-0.2\height}{\faGithub}~\href{https://github.com/usenameFD}{\underline{usenameFD}} \hspace{1em}
\raisebox{-0.2\height}{\faGlobe}~\href{https://usenamefd.github.io/}{\underline{usenamefd.github.io}}
\vspace{-8pt}
}

%  Do not put content here — only packages, formatting and macros.
% =====================================================================
\documentclass[letterpaper,11pt]{article}

\usepackage{latexsym}
\usepackage[empty]{fullpage}
\usepackage{titlesec}
\usepackage{marvosym}
\usepackage[usenames,dvipsnames]{color}
\usepackage{verbatim}
\usepackage{enumitem}
\usepackage[hidelinks]{hyperref}
\usepackage{fancyhdr}
\usepackage[english]{babel}
\usepackage{tabularx}
\usepackage{fontawesome5}
\usepackage{multicol}
\usepackage{graphicx}
\setlength{\multicolsep}{-3.0pt}
\setlength{\columnsep}{-1pt}
\input{glyphtounicode}
\usepackage{paracol}

\pagestyle{fancy}
\fancyhf{}
\fancyfoot{}
\renewcommand{\headrulewidth}{0pt}
\renewcommand{\footrulewidth}{0pt}

\addtolength{\oddsidemargin}{-0.5in}
\addtolength{\evensidemargin}{-0.5in}
\addtolength{\textwidth}{0.9in}
\addtolength{\topmargin}{-.3in}
\addtolength{\textheight}{0.9in}
\hypersetup{
    colorlinks=true,
    linkcolor=blue,
    urlcolor=blue,
    filecolor=blue,
    pdfpagemode=FullScreen
}
\urlstyle{same}

\raggedbottom
\raggedright
\setlength{\tabcolsep}{0in}

\titleformat{\section}{
  \vspace{-4pt}\scshape\raggedright\large\bfseries
}{}{0em}{}[\color{black}\titlerule \vspace{-5pt}]

\pdfgentounicode=1

\newcommand{\resumeItem}[1]{
  \item\small{{#1 \vspace{-2pt}}}
}

\newcommand{\resumeSubheading}[4]{
  \vspace{-2pt}\item
    \begin{tabular*}{1.0\textwidth}[t]{l@{\extracolsep{\fill}}r}
      \textbf{#1} & \textbf{\small #2} \\
      \textit{\small#3} & \textit{\small #4} \\
    \end{tabular*}\vspace{-7pt}
}

\newcommand{\resumeProjectHeading}[2]{
    \item
    \begin{tabular*}{0.97\textwidth}{l@{\extracolsep{\fill}}r}
      \small#1 & \textbf{\small #2}\\
    \end{tabular*}\vspace{-7pt}
}

\renewcommand\labelitemi{$\vcenter{\hbox{\tiny$\bullet$}}$}
\renewcommand\labelitemii{$\vcenter{\hbox{\tiny$\bullet$}}$}

\newcommand{\resumeSubHeadingListStart}{\begin{itemize}[leftmargin=0.0in, label={}]}
\newcommand{\resumeSubHeadingListEnd}{\end{itemize}}
\newcommand{\resumeItemListStart}{\begin{itemize}}
\newcommand{\resumeItemListEnd}{\end{itemize}\vspace{-5pt}}

% ---------------------------------------------------------------------
% \cvheader{positioning line}{focus line}
% ---------------------------------------------------------------------
\newcommand{\cvheader}[2]{%
\begin{center}
    {\Huge \scshape Franck Deba} \\[0.4em]
    {\small \textit{#1}} \\[0.25em]
    {\small #2}
\end{center}
\vspace{0.5pt}
\small
\raisebox{-0.1\height}{\faPhone}~+33 (0)6 05 77 47 51 \hspace{1em}
\raisebox{-0.2\height}{\faEnvelope}~\href{mailto:debafranck@gmail.com}{\underline{debafranck@gmail.com}} \hspace{1em}
\raisebox{-0.2\height}{\faLinkedin}~\href{https://linkedin.com/in/deba-franck}{\underline{deba-franck}} \hspace{1em}
\raisebox{-0.2\height}{\faGithub}~\href{https://github.com/usenameFD}{\underline{usenameFD}} \hspace{1em}
\raisebox{-0.2\height}{\faGlobe}~\href{https://usenamefd.github.io/}{\underline{usenamefd.github.io}}
\vspace{-8pt}
}

%  Do not put content here — only packages, formatting and macros.
% =====================================================================
\documentclass[letterpaper,11pt]{article}

\usepackage{latexsym}
\usepackage[empty]{fullpage}
\usepackage{titlesec}
\usepackage{marvosym}
\usepackage[usenames,dvipsnames]{color}
\usepackage{verbatim}
\usepackage{enumitem}
\usepackage[hidelinks]{hyperref}
\usepackage{fancyhdr}
\usepackage[english]{babel}
\usepackage{tabularx}
\usepackage{fontawesome5}
\usepackage{multicol}
\usepackage{graphicx}
\setlength{\multicolsep}{-3.0pt}
\setlength{\columnsep}{-1pt}
\input{glyphtounicode}
\usepackage{paracol}

\pagestyle{fancy}
\fancyhf{}
\fancyfoot{}
\renewcommand{\headrulewidth}{0pt}
\renewcommand{\footrulewidth}{0pt}

\addtolength{\oddsidemargin}{-0.5in}
\addtolength{\evensidemargin}{-0.5in}
\addtolength{\textwidth}{0.9in}
\addtolength{\topmargin}{-.3in}
\addtolength{\textheight}{0.9in}
\hypersetup{
    colorlinks=true,
    linkcolor=blue,
    urlcolor=blue,
    filecolor=blue,
    pdfpagemode=FullScreen
}
\urlstyle{same}

\raggedbottom
\raggedright
\setlength{\tabcolsep}{0in}

\titleformat{\section}{
  \vspace{-4pt}\scshape\raggedright\large\bfseries
}{}{0em}{}[\color{black}\titlerule \vspace{-5pt}]

\pdfgentounicode=1

\newcommand{\resumeItem}[1]{
  \item\small{{#1 \vspace{-2pt}}}
}

\newcommand{\resumeSubheading}[4]{
  \vspace{-2pt}\item
    \begin{tabular*}{1.0\textwidth}[t]{l@{\extracolsep{\fill}}r}
      \textbf{#1} & \textbf{\small #2} \\
      \textit{\small#3} & \textit{\small #4} \\
    \end{tabular*}\vspace{-7pt}
}

\newcommand{\resumeProjectHeading}[2]{
    \item
    \begin{tabular*}{0.97\textwidth}{l@{\extracolsep{\fill}}r}
      \small#1 & \textbf{\small #2}\\
    \end{tabular*}\vspace{-7pt}
}

\renewcommand\labelitemi{$\vcenter{\hbox{\tiny$\bullet$}}$}
\renewcommand\labelitemii{$\vcenter{\hbox{\tiny$\bullet$}}$}

\newcommand{\resumeSubHeadingListStart}{\begin{itemize}[leftmargin=0.0in, label={}]}
\newcommand{\resumeSubHeadingListEnd}{\end{itemize}}
\newcommand{\resumeItemListStart}{\begin{itemize}}
\newcommand{\resumeItemListEnd}{\end{itemize}\vspace{-5pt}}

% ---------------------------------------------------------------------
% \cvheader{positioning line}{focus line}
% ---------------------------------------------------------------------
\newcommand{\cvheader}[2]{%
\begin{center}
    {\Huge \scshape Franck Deba} \\[0.4em]
    {\small \textit{#1}} \\[0.25em]
    {\small #2}
\end{center}
\vspace{0.5pt}
\small
\raisebox{-0.1\height}{\faPhone}~+33 (0)6 05 77 47 51 \hspace{1em}
\raisebox{-0.2\height}{\faEnvelope}~\href{mailto:debafranck@gmail.com}{\underline{debafranck@gmail.com}} \hspace{1em}
\raisebox{-0.2\height}{\faLinkedin}~\href{https://linkedin.com/in/deba-franck}{\underline{deba-franck}} \hspace{1em}
\raisebox{-0.2\height}{\faGithub}~\href{https://github.com/usenameFD}{\underline{usenameFD}} \hspace{1em}
\raisebox{-0.2\height}{\faGlobe}~\href{https://usenamefd.github.io/}{\underline{usenamefd.github.io}}
\vspace{-8pt}
}


\begin{document}

\cvheader{Quantitative Risk Analyst}{FRTB \& Default Risk Charge $\cdot$ IRB credit risk (PD, LGD, MoC) $\cdot$ VaR / ES \& backtesting $\cdot$ Stress testing $\cdot$ Model validation}

%-----------EXPERIENCE-----------
\section{Professional Experience}
\resumeSubHeadingListStart
  \resumeSubheading
    {\textbf{Quantitative Analyst (Intern)} $|$ \textbf{PwC}}{Apr -- Oct 2026}
    {Risk Line of Service}{Paris, France}
    \resumeItemListStart
      \resumeItem{Modelling of the Default Risk Charge (DRC) under the FRTB capital framework}
      \resumeItem{Latent multi-factor default model calibrated by three competing methods; Monte Carlo DRC at the 99.9\% one-year quantile with sectioning confidence intervals, as required by supervisory guidance}
      \resumeItem{Benchmarked the internal model against the CRR3 standardised approach (rating $\rightarrow$ CQS $\rightarrow$ risk weight)}
      \resumeItem{Demonstrated that a validation policy relying on the point DRC alone can stay silent on real calibration differences, and identified what it should monitor instead}
    \resumeItemListEnd

  \resumeSubheading
    {\textbf{Quantitative Risk Analyst (Intern)} $|$ \textbf{Cr\'edit Agricole}}{Apr -- Sep 2025}
    {Risk Modeling Department}{Paris, France}
    \resumeItemListStart
      \resumeItem{Developed climate stress-testing methodologies using time series and Monte Carlo methods}
      \resumeItem{Built an automated backtesting tool in Python (Dash) for model validation}
      \resumeItem{Migrated models from SAS to a Python-compatible platform}
    \resumeItemListEnd

  \resumeSubheading
    {\textbf{Data Science Intern} $|$ \textbf{Groupe VYV -- Harmonie Mutuelle}}{Jun -- Aug 2024}
    {Assurance Pr\'evoyance}{Paris, France}
    \resumeItemListStart
      \resumeItem{Risk profiling by k-means clustering; survival analysis (Cox model) for prolonged-treatment risk and expenditure estimation}
    \resumeItemListEnd
\resumeSubHeadingListEnd

%-----------REGULATORY MODELLING-----------
\section{Regulatory Modelling Projects}
\resumeSubHeadingListStart
  \resumeProjectHeading
    {\textbf{Credit Scoring --- LCL} $|$ \emph{Python, Logistic Regression, WOE}} {Nov 2024 -- Mar 2025}
    \resumeItemListStart
      \resumeItem{Credit scoring model for default probability estimation under the IRB regulatory framework}
      \resumeItem{Homogeneous Risk Classes built with logistic regression, WOE, Jenks discretisation and machine-learning approaches}
      \resumeItem{Statistical tests for grade heterogeneity and within-grade homogeneity}
    \resumeItemListEnd

  \resumeProjectHeading
    {\textbf{Margin of Conservatism --- EY Quantitative Team} $|$ \emph{Python, Bootstrap}} {Nov 2024 -- Mar 2025}
    \resumeItemListStart
      \resumeItem{Computed Margin of Conservatism type C for PD modelling under the credit risk framework}
      \resumeItem{Challenged the bootstrap methodology for model uncertainty estimation and explored alternatives}
    \resumeItemListEnd

  \resumeProjectHeading
    {\textbf{LGD Segmentation --- Cr\'edit Mutuel Ark\'ea} $|$ \emph{Random Forest, XGBoost}} {Jan 2023 -- Apr 2024}
    \resumeItemListStart
      \resumeItem{Identified risk drivers for LGD segmentation; built business-relevant and statistically robust segments}
      \resumeItem{Challenged number-weighted vs.\ outstanding-weighted LGD approaches under ECB requirements}
    \resumeItemListEnd
\resumeSubHeadingListEnd

%-----------RISK MEASUREMENT PROJECTS-----------
\section{Risk Measurement Projects}
\resumeSubHeadingListStart
  \resumeProjectHeading
    {\href{https://usenamefd.github.io/projects/value-at-risk.html}{\textbf{VaR, Expected Shortfall and Backtesting}} $|$ \emph{Python, R}} {2025}
    \resumeItemListStart
      \resumeItem{Four families of estimators: historical (with bootstrap confidence intervals), parametric (Gaussian, EWMA, skewed-Student by maximum likelihood), Monte Carlo, and extreme value theory}
      \resumeItem{EVT by block maxima (GEV vs.\ Gumbel with likelihood-ratio testing) and peaks-over-threshold (GPD) with automatic threshold calibration and KS validation}
      \resumeItem{AR(1)--GARCH(1,1) filtering with GPD on standardised residuals to address volatility clustering}
      \resumeItem{Kupiec unconditional-coverage backtesting and adaptive rolling recalibration}
    \resumeItemListEnd

  \resumeProjectHeading
    {\href{https://usenamefd.github.io/projects/creditvar-copules.html}{\textbf{CreditVaR of a Receivables Portfolio}} $|$ \emph{Python, Monte Carlo}} {2025}
    \resumeItemListStart
      \resumeItem{Market-implied default probabilities, estimated recovery-rate distributions and copula-based dependence between risk factors}
      \resumeItem{Dependogram evidence of upper-tail dependence that a Gaussian copula cannot reproduce --- the mechanism behind the underestimation of joint defaults}
      \resumeItem{Monte Carlo CreditVaR at 99\%}
    \resumeItemListEnd

  \resumeProjectHeading
    {\href{https://usenamefd.github.io/projects/risques-asset-management.html}{\textbf{Liquidity, Rate and Counterparty Risk of a Portfolio}} $|$ \emph{Python}} {2025}
    \resumeItemListStart
      \resumeItem{Liquidity run-off profiles under normal and stressed conditions, with and without behavioural deformation; derived an optimal portfolio size}
      \resumeItem{Bond VaR by sensitivity (Z-VaR) vs.\ full repricing, and joint rate/credit-spread VaR via the credit triangle}
    \resumeItemListEnd

  \resumeProjectHeading
    {\href{https://usenamefd.github.io/projects/risque-climatique-communes.html}{\textbf{Physical Climate Risk Exposure Index}} $|$ \emph{Python, GeoPandas}} {2025}
    \resumeItemListStart
      \resumeItem{Built a municipality-level exposure index combining recognised natural-disaster hazards and population density, discretised by Jenks natural breaks and mapped}
      \resumeItem{Cross-referenced with physical risk indicators from regionalised climate projections}
    \resumeItemListEnd
\resumeSubHeadingListEnd

%-----------EDUCATION-----------
\section{Education}
\resumeSubHeadingListStart
  \resumeSubheading
    {Master M2MO --- Quantitative Finance and Data Science}{Sep 2025 -- Present}
    {Universit\'e Paris Cit\'e}{Paris, France}
  \vspace{3pt}
  \resumeSubheading
    {MSc in Data Science and Risk Management}{Aug 2023 -- Dec 2025}
    {ENSAI}{Rennes, France}
  \vspace{3pt}
  \resumeSubheading
    {MSc in Statistics applied to Economics}{Sep 2021 -- Aug 2023}
    {ISSEA -- CEMAC}{Yaound\'e, Cameroon}
  \vspace{3pt}
  \resumeSubheading
    {Bachelor's Degree in Mathematics}{Sep 2017 -- Aug 2020}
    {Higher Teacher Training College \& University of Yaound\'e I}{Yaound\'e, Cameroon}
\resumeSubHeadingListEnd

%-----------SKILLS-----------
\section{Technical Skills}
\begin{itemize}[leftmargin=0.15in, label={}]
  \small{\item{
    \textbf{Regulatory:} FRTB / CRR3 (DRC, standardised approach), Basel III, IRB (PD, LGD, Margin of Conservatism), ECB guidelines, climate stress testing \\
    \textbf{Methods:} VaR / Expected Shortfall, extreme value theory (GEV, POT--GPD), GARCH, copulas, Monte Carlo, bootstrap, logistic regression, WOE, Random Forest, XGBoost, survival analysis \\
    \textbf{Languages:} Python (NumPy, SciPy, pandas, scikit-learn, statsmodels), R, SAS, SQL, C++ (basics) \\
    \textbf{Tools:} Dash, Git, Linux, \LaTeX
  }}
\end{itemize}

%-----------CERTIFICATIONS-----------
\section{Certifications, Awards \& Languages}
\begin{itemize}[leftmargin=0.15in, label={}]
  \small{\item{
    \textbf{Certifications:} AMF (Autorit\'e des March\'es Financiers), Financial Markets (Yale), AI in Finance (World Bank) \\
    \textbf{Awards:} National Essay Competition Prize (2019), PK Fokam Prize for Science Students (2017) \\
    \textbf{Languages:} French (native), English (advanced, B2 --- TOEIC/TOEFL)
  }}
\end{itemize}

\end{document}
